\uuid{K6AQ}
\exo7id{2745}
\titre{exo7 2745}
\auteur{tumpach}
\organisation{exo7}
\datecreate{2009-10-25}
\isIndication{false}
\isCorrection{false}
\chapitre{Application linéaire}
\sousChapitre{Image et noyau, théorème du rang}
\cpge{mpsi-mp2i.matrices-applications-lineaires.b, pcsi.matrices-applications-lineaires.b, ptsi.matrices-applications-lineaires.b, mpsi-mp2i.matrices-applications-lineaires.c, pcsi.matrices-applications-lineaires.d, ptsi.matrices-applications-lineaires.d, bcpst1.applications-lineaires-matrices, bcpst2.applications-lineaires-matrices.c}
\module{Algèbre}
\niveau{L1}
\difficulte{}

\contenu{
\texte{
Soit la matrice $$B = \left( \begin{array}{cccc}1 & 2 & 3 & 1\\-1 & 2 & -1 &-3\\ -3 & 5 &  2 &-3\end{array}\right).$$
}
\begin{enumerate}
    \item \question{D\'eterminer une base du noyau de $B$.}
    \item \question{D\'eterminer une base de l'image de $B$.}
\end{enumerate}
}
